% !TEX root = ../../../main.tex
% Appendix entry: detection.threshold -- "Threshold to spans" (Appendix A).
\entryfield{Algorithm}
\code{_run} in \code{Adapters/detection_threshold.py}, for scores $s_0,\dots,s_{n-1}$, threshold
$\theta$, merge gap $g$ seconds and floor $d$ seconds:
\begin{enumerate}
\item The mask $b_i = (s_i > \theta)$: strictly greater. A NaN compares false, so NaN is never
  above the cut; $+\infty$ compares true.
\item \code{spans_from_mask}: pad the mask with false at both ends and take the maximal runs of
  true as half-open intervals $[a_j, e_j)$. Each span's score is $\max_{a_j \le i < e_j} s_i$.
\item \code{tidy_spans} (shared with \blockref{detection.signal_threshold}): convert the two
  settings to samples at the \texttt{Scores}' own rate, $G = \mathrm{round}(g f_s)$ and
  $D = \mathrm{round}(d f_s)$. Walking the spans in order, a span whose start lies fewer than $G$
  samples after the end of the previous (merged) span is joined to it: the merged span runs from
  the first start to the last end and scores the maximum of its parts. Then every span with
  $e - a < D$ is dropped. Merging comes first, so a fragmented event can be long enough to keep.
\end{enumerate}
At 0 both settings change nothing: no gap is less than 0 samples, and no span is shorter than 0.

\entryfield{Assumptions}
\begin{itemize}
\item \emph{Direction.} Only values above the threshold are kept; there is no ``below'' setting.
  To keep the low values of a score, the score itself must be turned over upstream.
\item \emph{Units.} The threshold is in the units of the incoming score, whatever they are; the
  block knows nothing about them. For \blockref{detection.matrix_profile} the scores are
  z-normalised distances, which run from 0 to $2\sqrt{m}$ and sit near $\sqrt{2m}$ for unrelated
  windows, so one cut means different things at different window lengths. For
  \blockref{detection.fluss} the score runs from 0 to 1, but its peaks rise with the window
  (that block's description).
\item \emph{Alignment.} The spans index the scores. Where a score belongs to the window
  \emph{starting} at a sample (matrix profile, FLUSS), a span $[a, e)$ is a run of window starts,
  and the signal those windows cover is $[a, e + m - 1)$. The block does not know $m$ and does not
  widen the span.
\item \emph{Missing values.} NaN is below every cut, so the NaN tail of a start-aligned profile and
  FLUSS's blanked ends never become spans, and a NaN inside a run splits it in two (merging can
  join it again).
\end{itemize}

\entryfield{Outputs}
A \texttt{SpanSet} of half-open sample intervals, counted from the start of the span, with one
score each (the maximum inside) and no labels. The executor writes each span to the
\code{detections} table. The block puts nothing in \code{meta}.

\entryfield{Cost}
One pass over the scores and one over the spans: $O(n)$ time and memory.

\entryfield{Failure modes}
\begin{itemize}
\item \emph{Gaps read as oddities.} \blockref{detection.matrix_profile} gives windows touching a
  NaN an infinite distance, not NaN, and an infinite score is above any cut: every NaN gap in the
  recording becomes a span of ``unusual'' windows, $m$ samples wider than the gap.
\item \emph{Discords, not motifs.} On a matrix profile the block marks the windows with no close
  look-alike. That is what the \texttt{mp\_threshold} template does (cut 8.0 at a 1-minute window);
  for repeated shapes use \blockref{detection.mp_motifs}.
\item \emph{Window offset.} Spans on start-aligned scores stop up to $m-1$ samples before the end
  of the signal that produced them, and start where the first high-scoring window starts, which
  may be before the event (above). This is the window-versus-event offset found with the older
  report tool (\code{Pipelines/motif_report/}); this block shares it.
\item \emph{A cut carried across settings.} A threshold tuned at one window length, one recording
  or one upstream block does not transfer; the \texttt{fluss\_regimes} template's 0.46 is for a
  300\,s window only.
\item \emph{Grazing.} With both tidying settings at 0, a score that crosses the cut for one sample
  is a one-sample span, and a noisy score near the cut becomes many short spans.
\end{itemize}

\entryfield{Validation}
\code{tests/test_adapter_threshold.py} checks the rule exactly on small hand-made score arrays:
every run past the cut is kept at the defaults, the floor drops short spans, seconds round to
samples at $f_s$, merging joins close spans and scores their maximum, and merging comes before the
floor. That is the whole of what the block does, so these tests check the answer, not only the
shape. Whether a given cut is a good one depends on the upstream block and is not tested here.

\entryfield{References}
No published method: a plain level crossing.

\entryfield{Tests}
\code{tests/test_adapter_threshold.py}; in chains, \code{tests/test_execution.py}
(\code{test_matrix_profile_to_threshold_chain_writes_detections}, which also shows the
start-aligned end: all windows of a 200-sample span with $m = 6$ give one span ending at 195),
\code{tests/test_manifest.py}, \code{tests/test_template_fluss_regimes.py} and
\code{tests/test_adapter_fluss.py}.
